\documentclass[tikz,border=8pt]{standalone}
\usepackage{tikz}
\usepackage{amsmath}
\usepackage{amssymb}
\usepackage{pgfplots}
\pgfplotsset{compat=1.18}
\usepackage{ctex}
\usetikzlibrary{calc,positioning,arrows.meta,matrix,trees}

% LC 200 岛屿数量 DFS 连通分量
% 4x5 网格 [[1,1,0,0,1],[1,1,0,0,0],[0,0,1,0,0],[0,0,0,1,1]]
% 4 个岛屿用不同颜色标注

\begin{document}
\begin{tikzpicture}[
  every node/.style={font=\small},
  sea/.style={draw=gray!50, rectangle, minimum size=0.85cm, thick, fill=blue!8},
  island1/.style={draw=green!60!black, rectangle, minimum size=0.85cm, very thick, fill=green!25},
  island2/.style={draw=orange!70!black, rectangle, minimum size=0.85cm, very thick, fill=orange!20},
  island3/.style={draw=purple!60!black, rectangle, minimum size=0.85cm, very thick, fill=purple!15},
  island4/.style={draw=red!60!black, rectangle, minimum size=0.85cm, very thick, fill=red!15},
]

\def\cs{1.0}

%% 标题
\node[font=\large\bfseries] at (3.5, 1.5)
  {岛屿数量（LC 200 Number of Islands）};
\node[font=\small,gray] at (3.5, 0.9)
  {DFS 标记连通分量，4 色标注 4 个岛屿，count = 4};

%% 4行5列网格（行0上，列0左）
%% 网格值：[[1,1,0,0,1],[1,1,0,0,0],[0,0,1,0,0],[0,0,0,1,1]]
% row0: 1 1 0 0 1
\node[island1] at (0*\cs, 0)       {\textbf{1}};
\node[island1] at (1*\cs, 0)       {\textbf{1}};
\node[sea]     at (2*\cs, 0)       {0};
\node[sea]     at (3*\cs, 0)       {0};
\node[island2] at (4*\cs, 0)       {\textbf{1}};

% row1: 1 1 0 0 0
\node[island1] at (0*\cs, -1*\cs)  {\textbf{1}};
\node[island1] at (1*\cs, -1*\cs)  {\textbf{1}};
\node[sea]     at (2*\cs, -1*\cs)  {0};
\node[sea]     at (3*\cs, -1*\cs)  {0};
\node[sea]     at (4*\cs, -1*\cs)  {0};

% row2: 0 0 1 0 0
\node[sea]     at (0*\cs, -2*\cs)  {0};
\node[sea]     at (1*\cs, -2*\cs)  {0};
\node[island3] at (2*\cs, -2*\cs)  {\textbf{1}};
\node[sea]     at (3*\cs, -2*\cs)  {0};
\node[sea]     at (4*\cs, -2*\cs)  {0};

% row3: 0 0 0 1 1
\node[sea]     at (0*\cs, -3*\cs)  {0};
\node[sea]     at (1*\cs, -3*\cs)  {0};
\node[sea]     at (2*\cs, -3*\cs)  {0};
\node[island4] at (3*\cs, -3*\cs)  {\textbf{1}};
\node[island4] at (4*\cs, -3*\cs)  {\textbf{1}};

%% 行列标注
\foreach \c in {0,...,4} {
  \node[font=\tiny,gray] at (\c*\cs, 0.58) {c\c};
}
\foreach \r in {0,...,3} {
  \node[font=\tiny,gray,anchor=east] at (-0.5, -\r*\cs) {r\r};
}

%% DFS 遍历顺序箭头（岛屿1内部）
\draw[->,>=Stealth,green!60!black,thick]
  (0*\cs+0.4, 0) -- (1*\cs-0.4, 0);
\draw[->,>=Stealth,green!60!black,thick]
  (0*\cs, -0.4) -- (0*\cs, -1*\cs+0.4);
\draw[->,>=Stealth,green!60!black,thick]
  (0*\cs+0.4, -1*\cs) -- (1*\cs-0.4, -1*\cs);

%% 岛屿标注（放在图右侧，不重叠）
\node[draw=green!60!black,fill=green!8,rounded corners=2pt,
      font=\scriptsize,inner sep=4pt]
  at (6.6, -0.4) {岛屿 1（绿色）\\(0,0)(0,1)(1,0)(1,1)};
\node[draw=orange!70!black,fill=orange!8,rounded corners=2pt,
      font=\scriptsize,inner sep=4pt]
  at (6.6, -1.4) {岛屿 2（橙色）\\(0,4)};
\node[draw=purple!60!black,fill=purple!8,rounded corners=2pt,
      font=\scriptsize,inner sep=4pt]
  at (6.6, -2.3) {岛屿 3（紫色）\\(2,2)};
\node[draw=red!60!black,fill=red!8,rounded corners=2pt,
      font=\scriptsize,inner sep=4pt]
  at (6.6, -3.2) {岛屿 4（红色）\\(3,3)(3,4)};

%% count 结果
\node[draw=gray!60,fill=yellow!15,very thick,rounded corners=4pt,
      font=\large\bfseries]
  at (3.5, -4.8) {count = 4};

%% DFS 说明
\node[draw=gray!50,fill=white,rounded corners=3pt,font=\small,
      inner sep=8pt,align=left]
  at (4.5, -6.5)
  {\textbf{DFS 连通分量 O(m×n) 时间：}\\[4pt]
   遍历每个格子，若值=1 且未访问：\\[2pt]
   \quad 计数器 count++，以该格为根做 DFS\\[2pt]
   \quad DFS 将所有四向相连的 1 标记为"已访问"\\[2pt]
   遍历结束，count 即为岛屿数量};

\end{tikzpicture}
\end{document}
